%% ═══════════════════════════════════════════════════════════════════════════
%% CONVENCAO DE COR NESTE ARQUIVO
%%   AZUL  (\corr{})  = alterado nesta rodada de correcoes (parecer LC)
%%   PRETO (\revisao{})         = rodada anterior, ja incorporada (sem cor)
%%   LARANJA (\pendente{})      = pendencia que exige acao de Marcela antes da entrega
%% ═══════════════════════════════════════════════════════════════════════════

\chapter[Produção Científica Derivada desta Dissertação]{\corr{Produção Científica Derivada desta Dissertação}}
\label{apend:producoes}

%% ALTERADO (LC, ponto 1): titulo trocado de "Producoes Cientificas em Andamento".
%% O novo titulo cobre os tres desfechos possiveis ate a entrega (aceito, recusado,
%% em avaliacao) sem precisar de nova alteracao. Se o artigo for ACEITO ate outubro,
%% atualizar apenas o campo Status abaixo.

\artigo{O artigo científico derivado desta pesquisa foi aceito para publicação:}
\begin{itemize}
  \item \textbf{Título:} Formal Verification of Critical Modules in the PX4
        Autopilot using Bounded Model Checking
  \item \textbf{Autores:} Marcela Alves Rodrigues, Lucas Carvalho Cordeiro,
        José Reginaldo Hughes de Carvalho
  \item \textbf{Veículo:} XVI Simpósio Brasileiro de Engenharia de Sistemas
        Computacionais (SBESC~2026), trilha de artigo completo, Salvador/BA,
        24 a 27 de novembro de 2026
  %% [12/10/2026] Status confirmado pela autora: aceito; a versao final e a
  %% enviada a esta revisao. Acrescentar paginas/DOI dos anais quando publicados.
  \item \textbf{Status:} \artigo{Aceito (submissão \#33507); versão final
        enviada em outubro de 2026}
  \item \textbf{Artefato:} \textit{harnesses}, logs de execução e instruções de
        reprodução arquivados no Zenodo sob o DOI~\artigo{\texttt{10.5281/zenodo.23282236}
        (versão \texttt{v1.2-sbesc}; DOI de conceito \texttt{10.5281/zenodo.21401958})} e
        disponíveis em \url{https://github.com/marcela2881/dissertacao-px4-esbmc}
\end{itemize}


% =============================================================================
\chapter[Harness Completo --- GPS uBlox]{\corr{Harness Completo --- GPS uBlox}}
\label{apend:harness-gps}
% =============================================================================

%% CAPITULO NOVO — atende ao ponto (2) [IMPORTANTE] do parecer do LC: o documento
%% reivindica reprodutibilidade e cita ~2.800 linhas de harness, mas nenhum harness
%% completo aparece no texto. Um harness integral resolve de uma vez as divergencias
%% de assinatura, a duvida sobre a pre-condicao e a verificacao das flags usadas.
%%
%% ACAO: colar o conteudo LITERAL de test_gps_driver.cpp no lugar indicado
%% (~/dissertacao-px4-esbmc/resultados/GPS/test_gps_driver.cpp).

\corr{Este apêndice reproduz na íntegra o \textit{harness} de
verificação do módulo GPS uBlox, \texttt{test\_gps\_driver.cpp}, desenvolvido
conforme a metodologia descrita no Capítulo~\ref{cap:metodologia}. Trata-se do
caso desenvolvido de ponta a ponta nesta dissertação: dele decorre o
contraexemplo apresentado na Seção~\ref{sec:contraexemplo-gps}}\artigo{, e
sobre ele foram feitas a verificação do \textit{patch} e a análise de
alcançabilidade discutidas no Capítulo~\ref{cap:experimentos}.}
%% [12/10/2026] "validacao em hardware" retirada: nao houve validacao em
%% hardware (ver Cap. 6, trabalhos futuros).

\corr{A reprodução integral, e não por fragmentos, permite ao leitor
verificar diretamente três pontos que fragmentos não sustentam: a assinatura
exata das funções sob verificação, as pré-condições impostas a cada teste, e
quais construções pertencem ao driver original do PX4 e quais foram introduzidas
pelo \textit{harness}.}

\corr{O arquivo define cinco funções de teste, selecionadas em
\texttt{main()} por uma variável não-determinística. Cada uma impõe pré-condições
distintas sobre o comprimento de entrada e sobre o estado inicial do buffer, o
que explica por que um mesmo módulo admite contraexemplos com valores diferentes
conforme o teste alcançado primeiro pelo verificador.}

%% [11/10/2026] PENDENCIA RESOLVIDA: harness literal incluido a partir de
%% codigo/test_gps_driver.cpp, copia byte a byte de
%% resultados/GPS/test_gps_driver.cpp (commit b535829), so com as quebras de
%% linha CRLF convertidas para LF. Acentos tratados pelo \lstset{literate}
%% de 000-main.tex.
\artigo{O arquivo tem 294 linhas (Tabela~\ref{tab:modulos}) e é reproduzido sem
edição a partir do repositório de artefatos. O contraexemplo da
Tabela~\ref{tab:resultados} vem do teste selecionado por
\texttt{test\_choice = 3}, \texttt{test\_gps\_real\_bit\_operation()}
(linhas 178--197).}

%% LC (resolvido 12/11/2026): as duas observacoes sobre o cabecalho e os
%% comentarios do harness viraram o paragrafo abaixo. O arquivo e reproduzido
%% literalmente, por isso os comentarios nao foram editados; a divergencia e
%% declarada no texto. Conferido contra gps.cpp e GpsDump.msg no commit 8be2590c15.
\corr{Como o arquivo é reproduzido sem edição, seus comentários internos são os
da versão de desenvolvimento e divergem, em três pontos, do que este texto
afirma. Prevalece o texto. (i)~O cabeçalho cita a versão~1.16 e a linha~643 de
\texttt{gps.cpp}; a verificação reportada refere-se ao commit
\texttt{8be2590c15}, em que \texttt{GPS::dumpGpsData()} começa na linha~727.
(ii)~Os comandos listados no rodapé do arquivo (\texttt{--unwind 8} e
\texttt{--unwind 5}) são anteriores à campanha; os resultados desta dissertação
foram obtidos com \texttt{--unwind 10 --no-unwinding-assertions
--overflow-check} (Seção~\ref{sec:config-experimental}). (iii)~Os comentários
dizem que a função foi extraída ``exatamente'' do driver, mas o \textit{harness}
é uma transcrição adaptada: a capacidade do buffer é 200~bytes, e não os 79 de
\texttt{uint8[79] data} em \texttt{GpsDump.msg}; a função recebe o buffer e o
modo ativo como parâmetros adicionais, em vez de usar os membros da classe; e o
campo \texttt{device\_id} da mensagem real não está na estrutura. A lógica do
laço, da guarda de espaço, do \texttt{memcpy} e da operação de bit é a do
driver. A diferença de capacidade é a mais relevante para a interpretação do
contraexemplo e é discutida na Seção~\ref{subsec:gps-inducao}.}

%% [12/10/2026] Conteudo de codigo/test_gps_driver.cpp colado aqui (antes era
%% \lstinputlisting): o apendice nao depende mais da pasta codigo/.
\begin{lstlisting}[
  caption={\artigo{\textit{Harness} \texttt{test\_gps\_driver.cpp} (literal)}},
  label={lst:harness-gps},
  language=C++,
  numbers=left,
  frame=single,
  basicstyle=\ttfamily\scriptsize,
  breaklines=true,
  extendedchars=true,
  literate={á}{{\'a}}1 {à}{{`a}}1 {â}{{\^a}}1 {ã}{{\~a}}1 {é}{{\'e}}1 {ê}{{\^e}}1
    {í}{{\'i}}1 {ó}{{\'o}}1 {ô}{{\^o}}1 {õ}{{\~o}}1 {ú}{{\'u}}1 {ç}{{\c c}}1
    {Á}{{\'A}}1 {À}{{`A}}1 {Â}{{\^A}}1 {Ã}{{\~A}}1 {É}{{\'E}}1 {Ê}{{\^E}}1
    {Í}{{\'I}}1 {Ó}{{\'O}}1 {Ô}{{\^O}}1 {Õ}{{\~O}}1 {Ú}{{\'U}}1 {Ç}{{\c C}}1]
/**
 * @file teste_gps_driver_real_esbmc.cpp
 * @author Dissertação Mestrado - Verificação Formal PX4 v1.16
 * 
 * OBJETIVO: Teste de código REAL extraído do PX4 GPS Driver
 * FUNÇÃO TESTADA: GPS::dumpGpsData() - linha ~643 de src/drivers/gps/gps.cpp
 * MÉTODO: Bounded Model Checking com ESBMC
 * ESTRUTURAS: Baseadas no USO real no código gps.cpp
 */

#include <assert.h>
#include <cstring>
#include <cstdint>

// ================== FUNÇÕES ESBMC ==================
extern int nondet_int();
extern uint8_t nondet_uint8();
extern size_t nondet_size_t();
extern bool nondet_bool();
extern void __ESBMC_assume(int condition);

// ================== ESTRUTURAS REAIS EXTRAÍDAS DO PX4 ==================
/**
 * ESTRUTURA REAL baseada no USO no gps.cpp:
 * - dump_data->data + dump_data->len
 * - dump_data->len += write_len  
 * - sizeof(dump_data->data)
 * - dump_data->len |= 1 << 7
 * - dump_data->timestamp = hrt_absolute_time()
 * - dump_data->instance = (uint8_t)_instance
 */

// Tamanho realista que funciona com uint8_t len
#define GPS_DUMP_DATA_SIZE 200

struct gps_dump_s {
    uint8_t data[GPS_DUMP_DATA_SIZE];    // Buffer principal
    uint8_t len;                         // Índice atual (usado com bitwise ops)
    uint8_t instance;                    // Instância GPS
    uint64_t timestamp;                  // hrt_absolute_time()
};

enum class gps_dump_comm_mode_t : int32_t {
    Disabled = 0,
    Full = 1,
    RTCM = 2
};

// ================== FUNÇÃO REAL EXTRAÍDA DO PX4 ==================

/**
 * FUNÇÃO ORIGINAL dumpGpsData() extraída EXATAMENTE do gps.cpp linha ~643
 * Lógica preservada: while loop, memcpy, aritmética de ponteiros, bit operations
 */
void dumpGpsData(uint8_t *data, size_t len, gps_dump_comm_mode_t mode, bool msg_to_gps_device, 
                 gps_dump_s *dump_data, gps_dump_comm_mode_t active_mode)
{
    // Verificação de modo (do código original)
    if (active_mode != mode || !dump_data) {
        return;
    }

    dump_data->instance = 0; // Simular: dump_data->instance = (uint8_t)_instance;

    // LOOP CRÍTICO REAL DO PX4
    while (len > 0) {
        size_t write_len = len;

        // CÁLCULO CRÍTICO: potencial underflow se dump_data->len > GPS_DUMP_DATA_SIZE
        if (write_len > GPS_DUMP_DATA_SIZE - dump_data->len) {
            write_len = GPS_DUMP_DATA_SIZE - dump_data->len;
        }

        // OPERAÇÃO CRÍTICA: memcpy com aritmética de ponteiros (do código real)
        memcpy(dump_data->data + dump_data->len, data, write_len);
        
        // ATUALIZAÇÕES (exatamente como no gps.cpp)
        data += write_len;
        dump_data->len += write_len;
        len -= write_len;

        // LÓGICA DE PUBLICAÇÃO (do código original)
        if (dump_data->len >= GPS_DUMP_DATA_SIZE) {
            // BIT OPERATION do código real
            if (msg_to_gps_device) {
                dump_data->len |= 1 << 7;  // Set bit 7
            }

            // Simular: dump_data->timestamp = hrt_absolute_time();
            dump_data->timestamp = 12345;
            
            // Simular: _dump_communication_pub.publish(*dump_data);
            // Reset para próxima iteração (do código real)
            dump_data->len = 0;
        }
    }
}

// ================== TESTES DE VERIFICAÇÃO FORMAL ==================

/**
 * TESTE 1: Verificar buffer bounds na operação memcpy real
 * PROPRIEDADE: memcpy nunca deve escrever além do buffer data[]
 */
void test_gps_real_buffer_bounds() {
    size_t input_len = nondet_size_t();
    bool msg_to_device = nondet_bool();
    
    // Entrada realista mas potencialmente perigosa
    __ESBMC_assume(input_len > 0 && input_len <= 300);
    
    uint8_t input_data[300];
    gps_dump_s dump_buffer;
    dump_buffer.len = nondet_uint8();
    __ESBMC_assume(dump_buffer.len < GPS_DUMP_DATA_SIZE);
    
    uint8_t initial_len = dump_buffer.len;
    
    // Chamar função REAL do PX4
    dumpGpsData(input_data, input_len, gps_dump_comm_mode_t::Full, msg_to_device, 
                &dump_buffer, gps_dump_comm_mode_t::Full);
    
    // PROPRIEDADE CRÍTICA: len nunca deve exceder tamanho do buffer
    // (considerando que pode ser resetado para 0)
    assert(dump_buffer.len <= GPS_DUMP_DATA_SIZE);
}

/**
 * TESTE 2: Verificar underflow no cálculo GPS_DUMP_DATA_SIZE - dump_data->len
 * PROPRIEDADE: Subtração deve ser segura mesmo se len estiver corrompido
 */
void test_gps_real_underflow_protection() {
    size_t input_len = nondet_size_t();
    
    __ESBMC_assume(input_len > 0 && input_len <= 50);
    
    uint8_t input_data[50];
    gps_dump_s dump_buffer;
    
    // CENÁRIO CRÍTICO: len próximo ou maior que GPS_DUMP_DATA_SIZE
    dump_buffer.len = nondet_uint8();
    __ESBMC_assume(dump_buffer.len >= GPS_DUMP_DATA_SIZE - 10);
    
    // Chamar função REAL
    dumpGpsData(input_data, input_len, gps_dump_comm_mode_t::Full, false, 
                &dump_buffer, gps_dump_comm_mode_t::Full);
    
    // PROPRIEDADE: Resultado deve ser coerente
    assert(dump_buffer.len <= GPS_DUMP_DATA_SIZE);
}

/**
 * TESTE 3: Verificar terminação do loop while(len > 0) real
 * PROPRIEDADE: Loop deve sempre terminar, mesmo com inputs extremos
 */
void test_gps_real_loop_termination() {
    size_t input_len = nondet_size_t();
    
    __ESBMC_assume(input_len > 0 && input_len <= 100);
    
    uint8_t input_data[100];
    gps_dump_s dump_buffer;
    dump_buffer.len = nondet_uint8();
    __ESBMC_assume(dump_buffer.len < GPS_DUMP_DATA_SIZE);
    
    // Chamar função REAL
    dumpGpsData(input_data, input_len, gps_dump_comm_mode_t::Full, false, 
                &dump_buffer, gps_dump_comm_mode_t::Full);
    
    // PROPRIEDADE: Se chegou aqui, loop terminou
    assert(true);
}

/**
 * TESTE 4: Verificar bit operation real do código
 * PROPRIEDADE: dump_data->len |= 1 << 7 deve ser segura
 */
void test_gps_real_bit_operation() {
    size_t input_len = nondet_size_t();
    
    __ESBMC_assume(input_len > 0 && input_len <= GPS_DUMP_DATA_SIZE + 10);
    
    uint8_t input_data[GPS_DUMP_DATA_SIZE + 10];
    gps_dump_s dump_buffer;
    dump_buffer.len = nondet_uint8();
    __ESBMC_assume(dump_buffer.len >= GPS_DUMP_DATA_SIZE - 5);
    
    uint8_t initial_len = dump_buffer.len;
    
    // Chamar função REAL que pode executar bit operation
    dumpGpsData(input_data, input_len, gps_dump_comm_mode_t::Full, true, 
                &dump_buffer, gps_dump_comm_mode_t::Full);
    
    // PROPRIEDADE: Operação bit deve ser válida OU len deve ser resetado
    assert(dump_buffer.len == 0 || dump_buffer.len <= GPS_DUMP_DATA_SIZE || 
           (dump_buffer.len & 0x7F) <= GPS_DUMP_DATA_SIZE);
}

/**
 * TESTE 5: Verificar edge case quando buffer está completamente cheio
 * PROPRIEDADE: Função deve lidar corretamente com buffer no limite
 */
void test_gps_real_full_buffer_edge_case() {
    size_t input_len = nondet_size_t();
    
    __ESBMC_assume(input_len > 0 && input_len <= 20);
    
    uint8_t input_data[20];
    gps_dump_s dump_buffer;
    
    // Buffer exatamente cheio
    dump_buffer.len = GPS_DUMP_DATA_SIZE;
    
    // Chamar função REAL
    dumpGpsData(input_data, input_len, gps_dump_comm_mode_t::Full, false, 
                &dump_buffer, gps_dump_comm_mode_t::Full);
    
    // PROPRIEDADE: Deve lidar graciosamente com buffer cheio
    assert(dump_buffer.len <= GPS_DUMP_DATA_SIZE);
}

// ================== MAIN PARA ESBMC ==================
int main() {
    int test_choice = nondet_int();
    __ESBMC_assume(test_choice >= 0 && test_choice < 5);
    
    switch(test_choice) {
        case 0:
            test_gps_real_buffer_bounds();
            break;
        case 1:
            test_gps_real_underflow_protection();
            break;
        case 2:
            test_gps_real_loop_termination();
            break;
        case 3:
            test_gps_real_bit_operation();
            break;
        case 4:
            test_gps_real_full_buffer_edge_case();
            break;
    }
    
    return 0;
}

/* 
 * ================================================================
 * DOCUMENTAÇÃO
 * ================================================================
 * 
 * METODOLOGIA DEMONSTRADA:
 * 
 * 1. EXTRAÇÃO DE CÓDIGO 100% REAL:
 *    - Função dumpGpsData() copiada exatamente de src/drivers/gps/gps.cpp
 *    - Estrutura gps_dump_s baseada no USO real observado no código
 *    - Lógica preservada: while loop, memcpy, bit operations, reset
 * 
 * 2. ESTRUTURA DERIVADA DO CÓDIGO REAL:
 *    - data[]: deduzido de sizeof(dump_data->data) 
 *    - len: deduzido de dump_data->len += write_len e operações bit
 *    - instance: observado em dump_data->instance = (uint8_t)_instance
 *    - timestamp: observado em dump_data->timestamp = hrt_absolute_time()
 * 
 * 3. PROPRIEDADES BASEADAS EM VULNERABILIDADES REAIS:
 *    - Buffer overflow em memcpy + aritmética de ponteiros
 *    - Integer underflow em GPS_DUMP_DATA_SIZE - dump_data->len
 *    - Loop infinito se write_len não decrementar len corretamente
 *    - Bit operation safety em len |= 1 << 7
 * 
 * 4. TÉCNICA DE VERIFICAÇÃO:
 *    - Bounded Model Checking com ESBMC
 *    - Entrada simbólica para tamanhos e estados do buffer
 *    - Assumptions para cenários críticos específicos
 *    - Assertions para propriedades de memory safety
 * 
 * COMANDOS DE EXECUÇÃO:
 * esbmc teste_gps_driver_real_esbmc.cpp --unwind 8 --timeout 300s
 * esbmc teste_gps_driver_real_esbmc.cpp --overflow-check --unwind 5
 * 
 * VULNERABILIDADES ALVO:
 * - Buffer overflow se dump_data->len corrompido
 * - Integer underflow na subtração GPS_DUMP_DATA_SIZE - len
 * - Aritmética de ponteiros insegura em memcpy
 * - Loop infinito se write_len == 0
 * 
 * DIFERENCIAL DESTA VERSÃO:
 * - Estruturas baseadas em USO real, não especulação
 * - Função extraída palavra por palavra do gps.cpp
 * - Tamanhos realistas compatíveis com uint8_t
 * - Testes focados em padrões observados no código real
 * 
 * ================================================================
 */
\end{lstlisting}


% =============================================================================
\chapter{Contraexemplos Gerados pelo ESBMC}
\label{apend:contraexemplos}
% =============================================================================

%% Prof. Lucas (qualificacao): "Poderias mostrar o contraexemplo gerado pelo ESBMC
%% no apendice." -- ATENDIDO: apendice adicionado.

%% ━━━━━━━━━━━━━━━━━━━━━━━━━━━━━━━━━━━━━━━━━━━━━━━━━━━━━━━━━━━━━━━━━━━━━━━
%% [RESOLVIDO] PENDENCIA DOS LOGS LITERAIS
%%
%% O parecer do LC apontou que os numeros deste apendice contradiziam os do Cap. 5
%% e que o Contraexemplo 1 era aritmeticamente inconsistente. A causa foi
%% identificada: o apendice MISTURAVA DOIS RUNS DISTINTOS do ESBMC v7.10.0.
%%
%%   Run A (relatorio interno de 25/01/2026):  temp_msb=64, temp_lsb=0
%%                                             -> Temp_uint11=512 -> 87,00 C
%%   Run B (figura do artigo, Issue #26865):   temp_msb=63, temp_lsb=95
%%                                             -> Temp_uint11=506 -> 86,25 C
%%
%% Ambos sao validos: temp_msb e temp_lsb sao totalmente nao-deterministicos, logo
%% qualquer par que ultrapasse 85 C satisfaz a violacao, e a escolha da testemunha
%% e do solver. O apendice havia ficado com a ENTRADA do run A e o RESULTADO do
%% run B — dai a inconsistencia (512 x 0,125 + 23 = 87,00, nao 86,25).
%%
%% CORRECAO APLICADA: adotado integralmente o run B, que e o publicado no artigo
%% #33507, na Fig. ce-imu-temp e na Issue #26865. Valores conferidos:
%%   506 x 0,125 + 23 = 86,25 C            [OK]
%%   506 = 0000000111111010 (binario)      [OK, confere com a figura]
%% ━━━━━━━━━━━━━━━━━━━━━━━━━━━━━━━━━━━━━━━━━━━━━━━━━━━━━━━━━━━━━━━━━━━━━━━

Este apêndice apresenta os contraexemplos gerados pelo ESBMC nos módulos IMU
BMI088 e GPS uBlox. Conforme discutido no Capítulo~\ref{cap:experimentos}, apenas
o do GPS uBlox sobreviveu à triagem como violação alcançável no modelo; os do
BMI088 são reproduzidos aqui por completude documental, ilustrando a etapa de
triagem da metodologia. \corr{Os contraexemplos da segunda campanha (RoboClaw,
batmon e candidatos adicionais) não são reproduzidos aqui; os registros
completos estão no repositório de artefatos, nas pastas indicadas na
Seção~\ref{sec:segunda-campanha}.}

%% [11/10/2026] Tracos substituidos por trechos LITERAIS dos logs do ESBMC
%% 8.1.0 + Boolector 3.2.4 (os mesmos do artigo). As versoes anteriores eram
%% transcricoes do v7.10.0, com valores (206/234) e estados que nao constam de
%% nenhum log. Unica edicao: vetores de bytes longos elididos ("... omitidos
%% ...") e o caminho interno do ESBMC abreviado como <esbmc>.
\corr{Cada traço inclui os valores concretos atribuídos pelo solver
SMT, o rastro de execução até o estado de falha e a propriedade violada reportada
pelo verificador.} \artigo{Os traços abaixo são trechos literais dos logs do
ESBMC~v8.1.0 com Boolector~3.2.4 depositados no repositório de artefatos (os
mesmos que fundamentam a Tabela~\ref{tab:resultados} e as Listagens do artigo);
a única edição é a elisão de vetores de bytes longos, indicada no próprio
texto. Os contraexemplos foram originalmente reportados} \corr{à comunidade PX4 nas
Issues~\#26865 (IMU BMI088)\footnote{\url{https://github.com/PX4/PX4-Autopilot/issues/26865}}
e~\#26866 (GPS uBlox)\footnote{\url{https://github.com/PX4/PX4-Autopilot/issues/26866}}
\artigo{a partir de execuções anteriores com o ESBMC~v7.10.0; os tempos e os
valores aqui reportados vêm das execuções com o ESBMC~v8.1.0, conforme a nota
de rodapé do artigo.}}

\corr{Cabe uma observação metodológica sobre a leitura destes traços.
As entradas dos \textit{harnesses} são não-determinísticas, de modo que o solver
é livre para escolher qualquer testemunha que satisfaça a condição de violação.
Execuções distintas do mesmo \textit{harness}, sob versões ou solvers diferentes,
produzem valores concretos diferentes para a mesma propriedade violada, sem que
isso constitua divergência de resultado.} \artigo{Os valores reproduzidos
abaixo são os das execuções com o ESBMC~v8.1.0, idênticos nas três execuções
de cada módulo.}

% =============================================================================
\section{Contraexemplo 1 --- Temperatura fora da faixa operacional (BMI088)}
\label{sec:contraexemplo-bmi088-temp}
% =============================================================================

Violação da propriedade $\phi_{\text{temp}}$ (Equação~\ref{eq:prop_temp}): a
temperatura calculada de $86{,}25\,^{\circ}$C excede o limite superior de
$85\,^{\circ}$C de operação definido no \textit{datasheet} do
BMI088~\cite{Bosch2021}. \corr{A conversão especificada pelo
fabricante é $T = \texttt{Temp\_int11} \times 0{,}125 + 23$, de modo que o valor
\texttt{Temp\_uint11 = 506} produz exatamente
$506 \times 0{,}125 + 23 = 86{,}25\,^{\circ}$C.}

%% [11/10/2026] Listagem substituida pelo trecho literal de
%% resultados/IMU_BMI088/resultado_temperatura_8.1.log. A anterior dizia
%% "ESBMC v7.10.0" e "State 7 - line 51" (o log diz linha 42) e trazia um
%% comando que nao consta do log.
%% [12/10/2026] Conteudo de codigo/ce_temperatura_bmi088.txt colado aqui (antes era
%% \lstinputlisting): o apendice nao depende mais da pasta codigo/.
\begin{lstlisting}[
  caption={\corr{Contraexemplo ESBMC para a violação da faixa de temperatura no BMI088} \artigo{(trecho literal do log, ESBMC~v8.1.0)}},
  label={lst:contraex_bmi088_temp},
  numbers=left,
  frame=single,
  language=bash,
  basicstyle=\ttfamily\scriptsize,
  breaklines=true]
ESBMC version 8.1.0 64-bit x86_64 linux
Solving with solver Boolector 3.2.4
[Counterexample]

State 2 file resultados/IMU_BMI088/test_bmi088_imu.cpp line 120 column 5 function test_temperature_calculation thread 0
----------------------------------------------------
  temp_msb = 63 (00111111)

State 3 file resultados/IMU_BMI088/test_bmi088_imu.cpp line 121 column 5 function test_temperature_calculation thread 0
----------------------------------------------------
  temp_lsb = 95 (01011111)

State 7 file resultados/IMU_BMI088/test_bmi088_imu.cpp line 42 column 5 function updateTemperature thread 0
----------------------------------------------------
  Temp_uint11 = 506 (00000001 11111010)

State 8 file resultados/IMU_BMI088/test_bmi088_imu.cpp line 48 column 9 function updateTemperature thread 0
----------------------------------------------------
  Temp_int11 = 506 (00000001 11111010)

State 11 file resultados/IMU_BMI088/test_bmi088_imu.cpp line 51 column 5 function updateTemperature thread 0
----------------------------------------------------
  temperature = 8.625000e+1f (01000010 10101100 10000000 00000000)

State 12 file resultados/IMU_BMI088/test_bmi088_imu.cpp line 123 column 5 function test_temperature_calculation thread 0
----------------------------------------------------
  temperature = 8.625000e+1f (01000010 10101100 10000000 00000000)

State 13 file resultados/IMU_BMI088/test_bmi088_imu.cpp line 126 column 5 function test_temperature_calculation thread 0
----------------------------------------------------
Violated property:
  file resultados/IMU_BMI088/test_bmi088_imu.cpp line 126 column 5 function test_temperature_calculation
  assertion temperature >= -40.0f && temperature <= 85.0f
  temperature >= -40.000000f && temperature <= 85.000000f

VERIFICATION FAILED
\end{lstlisting}

\corr{A ausência de validação de faixa antes da aritmética de
compensação térmica caracteriza} \artigo{CWE-1284 (\textit{Improper Validation of
Specified Quantity in Input}), subclasse de CWE-20 (\textit{Improper Input
Validation})}\corr{.}
%% [12/10/2026] CWE-1284 mantida, como no artigo (pendencia encerrada).
\corr{
Conforme discutido na Seção~\ref{subsec:imu-temperatura}, esta observação foi
reportada à comunidade PX4 e recusada pelos mantenedores, que objetaram à
abordagem de saturação proposta no PR~\#27575; permanece registrada como nota de
projeto, sem consequência de segurança confirmada e sem remediação aceita.}

% =============================================================================
\section{Contraexemplo 2 --- Contador FIFO excede capacidade do buffer (BMI088)}
\label{sec:contraexemplo-bmi088-fifo}
% =============================================================================

%% ALTERADO: valor canonico do artigo (8.704; a versao anterior trazia 1.536) e
%% status de falso positivo explicitado.
%% [11/10/2026] 8.704 -> 9.088: valor dos tres logs IMU_BMI088_run*.log e do
%% artigo. "replay" -> conferencia contra o codigo de producao. Pendencia do
%% traco RESOLVIDA com o trecho literal de medianas_logs/IMU_BMI088_run1.log.

\corr{Violação da propriedade que limita o contador do FIFO ao tamanho
do buffer físico:} \artigo{\texttt{count = 9088}} \corr{bytes excede em mais de oito vezes o
limite de \texttt{FIFO\_SIZE = 1.024} bytes. \textbf{Este candidato foi refutado
como falso positivo} \artigo{na conferência contra o código do} driver de produção não
instrumentado (Seção~\ref{subsec:imu-fifo}): o driver liberado protege a leitura
em três níveis, e o valor} \artigo{$9.088 \geq 1.024$} \corr{dispara a primeira guarda de
reinicialização, de modo que nenhum acesso fora dos limites é alcançável. O
candidato surgiu porque o \textit{harness} omitia essas guardas, e foi capturado
pela etapa obrigatória de} \artigo{conferência contra o código de produção}\corr{.}

\artigo{O traço é o das três execuções da Tabela~\ref{tab:resultados} para o
IMU BMI088 (\texttt{test\_choice = 2}): é deste candidato, e não do de
temperatura, o tempo de 0,1\,s registrado na tabela.}

%% [12/10/2026] Conteudo de codigo/ce_fifo_bmi088.txt colado aqui (antes era
%% \lstinputlisting): o apendice nao depende mais da pasta codigo/.
\begin{lstlisting}[
  caption={\artigo{Contraexemplo ESBMC para o contador do FIFO no BMI088 (trecho literal do log, ESBMC~v8.1.0)}},
  label={lst:contraex_bmi088_fifo},
  numbers=left,
  frame=single,
  language=bash,
  basicstyle=\ttfamily\scriptsize,
  breaklines=true]
ESBMC version 8.1.0 64-bit x86_64 linux
Solving with solver Boolector 3.2.4
[Counterexample]

State 2 file resultados/IMU_BMI088/test_bmi088_imu.cpp line 245 column 5 function main thread 0
----------------------------------------------------
  test_choice = 2 (00000000 00000000 00000000 00000010)

State 4 file resultados/IMU_BMI088/test_bmi088_imu.cpp line 138 column 5 function test_fifo_count_calculation thread 0
----------------------------------------------------
  fifo_len_0 = 128 (10000000)

State 5 file resultados/IMU_BMI088/test_bmi088_imu.cpp line 139 column 5 function test_fifo_count_calculation thread 0
----------------------------------------------------
  fifo_len_1 = 227 (11100011)

State 6 file resultados/IMU_BMI088/test_bmi088_imu.cpp line 61 column 5 function fifoReadCount thread 0
----------------------------------------------------
  FIFO_LENGTH_1_MASKED = 35 (00100011)

State 8 file resultados/IMU_BMI088/test_bmi088_imu.cpp line 141 column 5 function test_fifo_count_calculation thread 0
----------------------------------------------------
  count = 9088 (00100011 10000000)

State 9 file resultados/IMU_BMI088/test_bmi088_imu.cpp line 144 column 5 function test_fifo_count_calculation thread 0
----------------------------------------------------
Violated property:
  file resultados/IMU_BMI088/test_bmi088_imu.cpp line 144 column 5 function test_fifo_count_calculation
  assertion count <= FIFO_SIZE
  (unsigned long int)count <= FIFO_SIZE

VERIFICATION FAILED
\end{lstlisting}

% =============================================================================
\section[Contraexemplo 3 --- Underflow e acesso fora dos limites (GPS uBlox)]{\corr{Contraexemplo 3 --- Underflow e acesso fora dos limites (GPS uBlox)}}
\label{sec:contraexemplo-gps}
% =============================================================================

%% ALTERADO (LC, ponto 4): as secoes "Contraexemplo 3" e "Contraexemplo 4" foram
%% FUNDIDAS. A figura do contraexemplo e o log mostram UMA unica propriedade
%% violada — a falha de dereferencia no memcpy. O ESBMC encerra na primeira
%% violacao, logo nao existem dois tracos independentes. A divisao anterior
%% sugeria duas verificacoes distintas, o que nao corresponde a saida da
%% ferramenta. Isso tambem resolve a contagem: tres contraexemplos, tres secoes.

Violação da propriedade $\phi_{\text{gps}}$ (Equação~\ref{eq:prop_gps}). A
verificação de espaço disponível em \texttt{dumpGpsData()} \artigo{subtrai o nível
de preenchimento da capacidade}:

\begin{lstlisting}[language=C,basicstyle=\ttfamily\scriptsize,frame=none,numbers=none]
if (write_len > GPS_DUMP_DATA_SIZE - dump_data->len)
\end{lstlisting}

%% [11/10/2026] Aritmetica do harness corrigida (artigo, Sec. V-D).
Quando \texttt{dump\_data->len} excede \texttt{GPS\_DUMP\_DATA\_SIZE}
($=200$), \artigo{a diferença, calculada em \texttt{int}, é negativa ($-7$ para
\texttt{len = 207}) e, convertida para \texttt{size\_t} na comparação com
\texttt{write\_len}, torna-se um valor sem sinal muito grande}, tornando a guarda permanentemente falsa. O \texttt{memcpy} subsequente
escreve além dos limites do buffer alocado.

\corr{O ESBMC reporta uma única propriedade violada: a falha de
\artigo{desreferência} no \texttt{memcpy}. O \textit{underflow} aritmético é a causa raiz,
e não uma segunda violação independente --- a ferramenta encerra na primeira
violação encontrada. A classificação CWE-191 (\textit{Integer Underflow})
$\rightarrow$~CWE-119 (\textit{Memory Buffer Bounds Violation}) reflete essa
relação de causa e efeito.}

%% [11/10/2026] Listagem substituida pelo trecho literal de
%% medianas_logs/GPS_run1.log (identico em run2 e run3): input_len = 10,
%% dump_buffer.len = 207. A versao anterior (206/234, estados 1-5, "200 - 234")
%% nao corresponde a nenhum log. O aviso antigo sobre as "duas versoes da
%% figura" fica obsoleto: a figura foi refeita a partir do log.
%% [12/10/2026] Conteudo de codigo/ce_gps_ublox.txt colado aqui (antes era
%% \lstinputlisting): o apendice nao depende mais da pasta codigo/.
\begin{lstlisting}[
  caption={\corr{Contraexemplo ESBMC para o underflow com acesso fora dos limites no GPS uBlox} \artigo{(trecho literal do log, ESBMC~v8.1.0)}},
  label={lst:contraex_gps},
  numbers=left,
  frame=single,
  language=bash,
  basicstyle=\ttfamily\scriptsize,
  breaklines=true]
ESBMC version 8.1.0 64-bit x86_64 linux
Solving with solver Boolector 3.2.4
[Counterexample]

State 1 file resultados/GPS/test_gps_driver.cpp line 224 column 5 function main thread 0
----------------------------------------------------
  test_choice = 3 (00000000 00000000 00000000 00000011)

State 3 file resultados/GPS/test_gps_driver.cpp line 179 column 5 function test_gps_real_bit_operation thread 0
----------------------------------------------------
  input_len = 10 (00000000 00000000 00000000 00000000 00000000 00000000 00000000 00001010)

State 6 file resultados/GPS/test_gps_driver.cpp line 183 column 5 function test_gps_real_bit_operation thread 0
----------------------------------------------------
  input_data = { ... 210 bytes omitidos ... }

State 7 file resultados/GPS/test_gps_driver.cpp line 184 column 5 function test_gps_real_bit_operation thread 0
----------------------------------------------------
  dump_buffer = { .data={ ... 200 bytes omitidos ... }, .len=0,
    .instance=0, .anon_pad$3=281474976710655, .timestamp=0xFFFFFFFFFFFFFFFF }

State 8 file resultados/GPS/test_gps_driver.cpp line 185 column 23 function test_gps_real_bit_operation thread 0
----------------------------------------------------
  dump_buffer.len = { .data={ ... 200 bytes omitidos ... }, .len=207,
    .instance=0, .anon_pad$3=281474976710655, .timestamp=0xFFFFFFFFFFFFFFFF }

State 21 file <esbmc>/c2goto/library/string.c line 280 column 5 function __memcpy_impl thread 0
----------------------------------------------------
Violated property:
  file <esbmc>/c2goto/library/string.c line 280 column 5 function __memcpy_impl
  dereference failure: Access to object out of bounds

VERIFICATION FAILED
\end{lstlisting}

\corr{O contraexemplo foi obtido com o solver Boolector~3.2.4.
Conforme registrado na Seção~\ref{sec:ameacas-validade}, esta é a única configuração que
descarrega a violação dentro do limite de tempo estabelecido, o que caracteriza o
achado como sensível ao solver --- uma das lições metodológicas desta pesquisa.}
\artigo{O estado violador (\texttt{len = 207}) é alcançável no modelo, mas não no
driver de produção, conforme a prova indutiva da
Seção~\ref{subsec:gps-inducao}.} \corr{O
\textit{harness} local \texttt{dump\_buffer} corresponde ao buffer
\texttt{dump\_data} repassado ao driver.}

Em plataformas Pixhawk sem unidade de gerenciamento de memória, uma escrita fora
dos limites pode corromper memória adjacente; sob o modelo de memória plana do
NuttX, não necessariamente gera uma falha limpa. O impacto preciso depende do
arranjo definido em tempo de ligação\artigo{, que não foi analisado nesta
dissertação; quanto à alcançabilidade, o estado violador não ocorre no driver de
produção (Seção~\ref{subsec:gps-inducao}).}


%% [12/10/2026] Apendice de VCCs (apend:vcc) RETIRADO: VCCs omitidas, como no
%% artigo. Os valores continuam nos logs do repositorio e em
%% ./tabela_ii_dos_logs.sh (Bar 93, Bat 49, EKF2 63, Expo 278, GPS 4463,
%% IMU 120, MAVLink 3059, Mag 50, Nav 50, PWM 38, RC 56, VTOL 110).
